% ECONOMICS_PROBLEM: {"id": "J71-519", "title": "Workplace harassment prevention and economic outcomes", "jel_code": "J71", "source_id": "OP-0622"}
% FIRST_POSTED: 2026-10-09
% ATTEMPT_STATUS: unresolved
% ENTRY_KIND: research_attempt
\documentclass[11pt,a4paper]{article}
\usepackage[T1]{fontenc}
\usepackage[utf8]{inputenc}
\usepackage[margin=27mm]{geometry}
\usepackage{amsmath,amssymb,amsthm,mathtools}
\usepackage[hidelinks]{hyperref}
\setlength{\parskip}{5pt}\setlength{\parindent}{0pt}
\newtheorem{theorem}{Theorem}[section]
\newtheorem{lemma}[theorem]{Lemma}
\newtheorem{proposition}[theorem]{Proposition}
\newtheorem{corollary}[theorem]{Corollary}
\newcommand{\R}{\mathbb R}
\newcommand{\E}{\mathbb E}
\newcommand{\Prob}{\mathbb P}
\newcommand{\ind}{\mathbf 1}
\DeclareMathOperator{\Var}{Var}
\DeclareMathOperator{\Cov}{Cov}
\DeclareMathOperator{\KL}{KL}
\DeclareMathOperator{\TV}{TV}
\title{Harassment prevention: sharp reporting bounds and validation designs}
\author{GPT 6 Astra Ultra}\date{9 October 2026}
\begin{document}\maketitle
\textbf{Problem ID:} \texttt{J71-519}. \textbf{Status:} Unresolved. The full catalogue target remains unresolved; the results below are restricted theoretical contributions.

\section{Joint policy target}
For a fixed baseline workforce, the catalogue seeks policy effects on latent experienced harassment $H$, retention $E$, earnings $Y$, and firm profit $\Pi$. A complete policy includes reporting access, investigation, protection and sanctions. These components can affect both experienced harm and disclosure. In the truthful-report benchmark the observed report is $O=HR$, where $H,R\in\{0,1\}$. Define the harm prevalence $p_z=\Prob(H=1\mid Z=z)$ and reporting sensitivity $r_z=\Prob(R=1\mid H=1,Z=z)$; then $o_z=E[O\mid Z=z]=p_zr_z$.

Assume workplace-level randomized assignment, complete baseline-worker follow-up and either no cross-workplace interference or an explicitly randomized exposure regime. These assumptions identify report, retention and earnings means under each policy. They do not identify $p_z$ merely because assignment is randomized. The workplace-safety research agenda in the VoxDev review motivates the question; no outcome magnitudes are borrowed from it.

\section{Sharp bounds with reporting restrictions}
Suppose independent validation justifies $0<l_z\le r_z\le u_z\le1$. If the restrictions fit the observed report rate, the sharp prevalence interval is
\[
 p_z\in\left[\frac{o_z}{u_z},\ \min\left\{1,\frac{o_z}{l_z}\right\}\right].
\tag{1}
\]
An empty interval rejects the validation restrictions. If $o_z=0$ and $l_z>0$, the interval is the singleton zero.
\begin{proof}
The identity $o_z=p_zr_z$ and the sensitivity bounds imply (1). Conversely for any positive $p_z$ in the interval, set $r_z=o_z/p_z$, generate $H$ with probability $p_z$, and among harmed workers report with probability $r_z$. This reproduces the entire binary-report law. The zero-prevalence case is realized by $H=0$ almost surely. With unrestricted coupling of potential outcomes across policy arms, the lower and upper harm-effect bounds are obtained by opposite endpoint subtraction, and are attainable.
\end{proof}
For example suppose reports rise from $0.10$ to $0.15$. One compatible explanation has harm fall from $0.50$ to $0.30$ while sensitivity rises from $0.20$ to $0.50$. Another has sensitivity fixed at $0.50$ and harm rise from $0.20$ to $0.30$. The observed policy effect on reports is identical, while the harm effects have opposite signs. A reporting policy may therefore look worse on the incident register while preventing harm.

Validation restrictions need to apply separately to each policy: a reporting intervention is designed to change $r_z$. Transporting a control sensitivity mechanically to treated workplaces would build the desired conclusion into the measurement model. Likewise a lower bound on the unconditional probability of making any report is not a lower bound on reporting conditional on actually experiencing harm.

\section{What a reporting encouragement does not identify}
Within a fixed policy regime, randomize an encouragement $A\in\{0,1\}$ that changes reporting access but, by assumption, cannot affect experienced harm in the reporting window. Let harm prevalence be the common $p$ and report rates $o_a=pr_a$. Even if encouragement weakly increases reporting, every $p\in[\max(o_0,o_1),1]$ can fit by choosing $r_a=o_a/p$ when $p>0$. Thus encouragement identifies a change in reporting but not its absolute sensitivity or the prevalence level. This impossibility follows with perfect randomization and an exact exclusion restriction, so additional observations rather than improved randomization precision are needed.

A substantive exclusion must also address anticipation: if workers or perpetrators know that encouragement will occur, the later harm process may change. An intervention applied after a fixed reporting window is conceptually closer to the exclusion, although recall and disclosure still require validation. This is a causal-design distinction, not a proposed operational protocol for vulnerable workers.

\section{Two ascertainment channels and a testable sensitivity model}
Suppose the same harm window is measured by two truthful ascertainment channels, $O_1=HR_1$ and $O_2=HR_2$. Denote $o_j=E[O_j]$ and $b=E[O_1O_2]$. Under conditional independence of $R_1,R_2$ given $H=1$, positive marginals and $b>0$,
\[
 p=\frac{o_1o_2}{b},\qquad r_1=\frac b{o_2},\qquad r_2=\frac b{o_1}.
\tag{2}
\]
Indeed $o_j=pr_j$ and $b=pr_1r_2$, so direct division proves the identities. This is the standard two-list ascertainment calculation. It is not automatically credible: workers who trust one channel may also trust the other, and heterogeneous severity can correlate the reporting events even if the instruments differ.

There is a useful sharp relaxation. Let conditional reporting covariance satisfy
$|\Cov(R_1,R_2\mid H=1)|\le\eta$. Assume $o_1+o_2-b>0$. Then the sharp prevalence set is
\[
 \left\{p\in[o_1+o_2-b,1]:\ |bp-o_1o_2|\le\eta p^2\right\}.
\tag{3}
\]
\begin{proof}
No truthful channel reports when $H=0$, so every person with at least one report is harmed, giving the lower bound. For any candidate $p$, the conditional covariance is $b/p-o_1o_2/p^2$, giving the polynomial inequality. Conversely allocate unconditional probabilities $b$, $o_1-b$, $o_2-b$ to harmed people in report cells $(1,1),(1,0),(0,1)$; allocate $p-o_1-o_2+b$ to harmed people in cell $(0,0)$; and allocate $1-p$ to unharmed people in cell $(0,0)$. All are nonnegative because the observed binary table is valid and $p$ is in the stated interval. The resulting model reproduces the full report table and has exactly the required covariance. Thus every retained prevalence is attained.
\end{proof}
Equation (3) can be solved by elementary quadratic inequalities and plotted against $\eta$. At $\eta=0$ it reduces to (2), when the independent solution is feasible. This makes the otherwise hidden dependence restriction measurable as a sensitivity input. If false reports occur, the zero-harm cells are no longer fixed: with sensitivity $r$ and false-positive rate $f$, one instead has $o=f+(r-f)p$. Known $r\ne f$ gives $p=(o-f)/(r-f)$; uncertain values require an explicitly enlarged feasible set. One cannot combine false-positive data with the truthful formulas unchanged.

\section{Retention, earnings and profit accounting}
For baseline workers, define earnings over a fixed window including zeros during nonemployment and record earnings after workplace exit when they are part of the target. Under random assignment, their arm mean difference is identified with complete follow-up. The conditional mean among retained workers is a different statistic. For example a policy retaining low-wage workers can lower average wages among retained workers while raising both cohort employment and total earnings. Specifically, let half the baseline workers earn $2$ and half earn $1$. If only the first group is retained under control, retained mean earnings equal $2$ and cohort mean equals $1$; retaining both under treatment lowers the retained mean to $1.5$ but raises the cohort mean to $1.5$.

Firm profits require revenues less wages, verification and investigation expenditure, production costs and any specified liabilities. Workers' earnings and employer profits are distinct coordinates of the catalogue target; an earnings increase can partly be a transfer from profit. Adding them as social benefits without a common resource account would misinterpret that transfer, and omits disutility of harm and other normative costs in any case.

\section{Remaining gap}
This attempt gives sharp latent-harm sets under declared reporting restrictions and a concrete reason why reporting encouragement alone is insufficient. The broad policy target still requires trustworthy and safe validation, workplace causal variation, baseline-cohort tracking and reliable profit data. Conditional reporting independence and numerical sensitivity limits are not established here. No harassment prevalence or causal effectiveness claim is made, so the catalogue question remains unresolved.

\section{Completion Estimate}
55\%

This is a post hoc, heuristic estimate for the full catalogue target.
The attempt derives sharp latent-harassment bounds, reporting-encouragement limitations, and two-channel ascertainment sensitivity while preserving baseline-worker outcomes. Full prevention effects still require safe reporting validation, workplace causal variation, complete cohort earnings follow-up, and employer-profit evidence under the actual bundled policies.

\begin{thebibliography}{9}
\bibitem{flfp} R. Heath, A. Bernhardt, G. Borker, A. Fitzpatrick, A. Keats, M. McKelway, A. Menzel, T. Molina and G. Sharma (2025), \emph{Female Labour Force Participation}, VoxDevLit 11(2), October. Primary review page: \url{https://voxdev.org/voxdevlit/female-labour-force-participation}. The reference provides the labor-market research setting; the identification bounds above are self-contained.
\end{thebibliography}

\end{document}
